%% kapa: exam=1 round=19 year=2008 period=1 form=A subject=civil_law subject_ko=민법 seq=35 number=35
%% source: Q-Net 감정평가사 기출문제 (cst00309) artlSeq=1001962 file=extracted/1008/1001962/2008년 제19회 감정평가사 1차시험 기출문제/1교시.hwp
%% conversion: kordoc 4.24.51 -> tools/gen1.py (auto); review=proofread-sonnet; flags=none
\begin{question}{35}
권리질권에 관한 설명으로 틀린 것은? (다툼이 있는 경우에는 판례에 의함)
\begin{choices}
  \item 지시채권의 입질은 그 증서에 배서하여 질권자에게 교부하여야 효력이 생긴다.
  \item 양도금지의 특약이 있는 채권은 질권의 대상으로 할 수 없지만, 양도금지의 특약으로 선의의 제3자에게는 대항할 수 없으므로, 질권자가 선의라면 유효하게 질권을 취득할 수 있다.
  \item 채권질권자가 질권의 목적인 채권을 ‘직접 청구할 수 있다’는 의미는 질권설정자의 이름으로 청구하고 추심할 권한을 갖는다는 것을 말한다.
  \item 무기명채권의 입질은 증서를 질권자에게 교부함으로써 그 효력이 생긴다.
  \item 은행이 대출채권의 담보로 자기에 대한 예금채권을 질권의 목적으로 하는 경우와 같이, 질권자 자신에 대한 채권도 질권의 대상이 된다.
\end{choices}
\end{question}
\answer{3}
